\begin{textsample}{NEG Dim 4 – 2012\_c5 – Score -1.00 – 2012\_c5/enade\_2012\_tec\_proc\_ger\_2c09315dc6.txt}  \label{ex:f4_neg_102}
CAPÍTULO 5 DISTRIBUIÇÃO DOS CONCEITOS 5.1 PANORAMA NACIONAL DA DISTRIBUIÇÃO DOS CONCEITOS A Tabela 5.1 apresenta a quantidade e distribuição de cursos de Tecnologia em Processos Gerenciais participantes do ENADE/2012, por faixa de conceito e Grande Região.
A diferença entre os cursos tabulados neste capítulo e no capítulo 2 são os cursos sem conceito, em princípio, aqueles sem alunos concluintes que participassem da prova ou que tivessem tão somente um aluno concluinte11.
Observando -se os dados da Tabela 5.1, nota -se que, dos 178 cursos participantes, 66 ( 37,1 \% ) classificaram -se com conceito 3, o valor modal.
Este foi também o conceito modal nas regiões Sudeste ( 37,1 \% ) e Sul ( 50,0 \% ).
Nas regiões Norte, Nordeste e Centro-Oeste o conceito modal foi 2, com respectivamente, 71,4 \%, 35,0 \% e 33,3 \%.
O conceito 2 foi o segundo mais frequente em nível nacional ( 30,9 \%, correspondendo a 55 cursos ) e o conceito 4, o terceiro ( 16,3 \%, correspondendo a 29 cursos ).
Houve, ainda, 12 cursos ( 6,7 \% ) que receberam conceito 5 e outros nove cursos ( 5,1 \% ) que receberam conceito 1.
Apenas sete dos 178 cursos de Tecnologia em Processos Gerenciais ficaram sem conceito ( SC ).
A região Norte participou com sete cursos ou 3,9 \% do total nacional.
Nesta região, como já comentado, o conceito 2 foi atribuído a cinco cursos dos sete participantes, o que equivale a 71,4 \% do total regional.
Nenhum curso ficou sem conceito ( SC ).
Os demais foram avaliados com conceitos 3 e 4 ( um curso cada, correspondendo a 14,3 \% ).
Nenhum curso foi \textbf{alocado} aos conceitos 1 ou 5. 11 Estes cursos com somente um concluinte foram considerados no capítulo 2.
A região Nordeste participou com 20 cursos ou 11,2 \% do total nacional.
Como já comentado, destes, sete cursos, 35,0 \% em termos regionais, obtiveram conceito 2, o conceito modal para a região.
Os demais receberam os conceitos 1 ( quatro cursos, 20,0 \% ), 3 ( cinco cursos, 25,0 \% ), 4 ( dois cursos, 10,0 \% ) e 5 ( também dois cursos ).
Nessa região, nenhum dos cursos ficou sem conceito ( SC ).
Dos 89 cursos participantes da região Sudeste, 33 ( 37,1 \% ) obtiveram conceito 3, o conceito modal.
O conceito 1 foi atribuído a três cursos ( 3,4 \% ) e o conceito 2, a 25 ( 28,1 \% ). 16 cursos ( 18,0 \% ) receberam o conceito 4 e sete outros ( 7,9 \% ) receberam o conceito 5.
Apenas cinco cursos ficaram sem conceito ( SC ).
A região Sul contou com 50 cursos ou 28,1 \% do total regional.
A predominância do conceito 3, como já comentado, foi de 50,0 \%, correspondentes a 25 dos 50 cursos participantes na região Sul.
O conceito 2 foi atribuído a 14 cursos ( 28,0 \% ), o conceito 4 a sete cursos ( 14,0 \% ) e o conceito 5, a três cursos ( 6,0 \% ).
Apenas um dos cursos da região Sul ficou sem conceito ( SC ), enquanto nenhum foi \textbf{alocado} ao conceito 1.
Exatamente 1/3 ( 33,3 \% correspondendo a quatro cursos ) dos 12 cursos participantes na região Centro-Oeste, como já comentado, recebeu conceito 2, o conceito modal.
Nesta região apenas um curso ficou sem conceito ( SC ).
Os demais cursos foram avaliados com conceito 1 ( dois cursos, 16,7 \% ), conceito 3 ( também dois cursos ) e conceito 4 ( três cursos, 25,0 \% ).
Nenhum curso foi \textbf{alocado} no conceito 5. 5.2 CONCEITOS POR CATEGORIA ADMINISTRATIVA E POR GRANDE REGIÃO A Tabela 5.2 apresenta a distribuição dos cursos participantes do ENADE/2012 de Tecnologia em Processos Gerenciais, por Categoria Administrativa, de acordo com os conceitos por eles alcançados, segundo as Grandes Regiões brasileiras.
Dos 178 cursos participantes, 26 ( 14,6 \% ) eram ministrados em instituições públicas e 152 ( 85,4 \% ), em privadas.
De acordo com as informações da Tabela 5.2, em termos nacionais, não houve equilíbrio na obtenção do conceito 5 entre instituições públicas e privadas : das 12 IES que receberam este conceito, cinco eram públicas e sete privadas.
Dos 26 cursos participantes de IES públicas, o conceito 4 foi o valor modal, atribuído a dez cursos ( 38,5 \% ).
Nesta categoria, apenas dois dos cursos ficaram sem conceito ( SC ).
Entre os demais cursos de instituições públicas participantes, um obteve conceito 1 ( 3,8 \% da categoria ) e oito receberam conceito 3 ( 30,8 \% ).
Como já comentado, os cinco restantes ( 19,2 \% ) receberam o conceito 5.
Na rede privada, o conceito modal foi 3, com 58 cursos dos 152 da categoria.
Nesta categoria administrativa, cinco dos cursos ficaram sem conceito ( SC ).
Entre os demais cursos participantes, oito receberam conceito 1 e 55, conceito 2.
O conceito 4 foi atribuído a 19 cursos e o conceito 5, como já comentado, a sete cursos.
Na análise por região, observa -se que, na região Norte, as instituições privadas concentraram todos os sete cursos participantes, dos quais nenhum ficou sem conceito ( SC ).
O conceito modal para as instituições privadas na região foi 2, com cinco cursos, correspondendo a 71,4 \%.
Dos dois cursos restante, um recebeu conceito 3 e o outro, conceito 4.
As instituições públicas não participaram com nenhum curso.
Na região Nordeste, a rede privada concentrou todos os 20 cursos participantes.
Nesta Categoria Administrativa na região, predominaram os cursos com conceito 2, sete cursos correspondendo a 35,0 \% dos cursos oferecidos por IES privada no Nordeste.
Os demais foram avaliados com conceitos 1 ( quatro cursos correspondendo a 20,0 \% ), 3 ( cinco cursos correspondendo a 25,0 \% ), 4 ( dois cursos correspondendo a 10,0 \% ) e 5 ( também dois cursos ).
As instituições públicas dessa região não participaram com nenhum curso.
Na região Sudeste, a proporção de cursos da rede privada foi de 80,9 \%, correspondendo a 72 dos 89 cursos participantes.
Nesta categoria, na região Sudeste, o conceito modal foi 3 ( 30 cursos ).
Nesta combinação de Categoria Administrativa e Grande Região, apenas quatro cursos não receberam conceito ( SC ).
Os demais foram avaliados com conceito 1 ( dois cursos ), conceito 2 ( 25 cursos ), conceito 4 ( oito cursos ) e conceito 5 ( três cursos ).
Entre os 17 cursos em instituições públicas na região Sudeste, a categoria modal foi a que recebeu conceito 4 ( oito cursos ).
Apenas um dos cursos ficou sem conceito ( SC ) e nenhum recebeu o conceito 2.
Os demais foram avaliados com os conceitos 1 ( um curso ), 3 ( três cursos ) e 5 ( quatro cursos)..
As instituições privadas concentraram 42 dos 50 cursos participantes da região Sul, 84,0 \% do total regional.
Desses, 20 ficaram com conceito 3, o conceito modal.
Os demais foram avaliados com conceito 2 ( 14 cursos ), conceito 4 ( seis cursos ) e conceito 5 ( dois cursos ).
Nesta combinação de Categoria Administrativa e Grande Região, nenhum curso ficou sem receber conceito ( SC ).
As instituições públicas na região Sul participaram com oito cursos ( 16,0 \% ), dos quais apenas um ficou sem conceito ( SC ).
O conceito modal também foi 3, com cinco cursos.
Além disso, os conceitos 4 e 5 foram recebidos por um curso cada.
Nenhum curso foi \textbf{alocado} aos conceitos 1 e 2.
Na região Centro-Oeste, 11 dos 12 cursos participantes eram de instituições privadas ( 91,7 \% em termos regionais ), dos quais apenas um curso ficou sem conceito ( SC ).
Dos demais cursos, dois receberam conceito 1, quatro receberam conceito 2, conceito modal, outros dois receberam conceito 3 e os dois cursos restantes, conceito 4.
Nenhum curso foi \textbf{alocado} ao conceito 5.
As instituições públicas participaram com somente um curso, o qual foi avaliado no conceito 4. 5.3 CONCEITOS POR ORGANIZAÇÃO ACADÊMICA E POR GRANDE REGIÃO Na Tabela 5.3 encontra -se a distribuição dos conceitos atribuídos aos cursos participantes do ENADE/2012 na Área de Tecnologia em Processos Gerenciais, por Organização Acadêmica, segundo as Grandes Regiões brasileiras.
Dos 178 cursos participantes, 56 eram oferecidos em Universidades, 26 em Centros Universitários e os demais 96 em Faculdades.
Esta distribuição corresponde a, respectivamente, 31,5 \%, 14,6 \% e 53,9 \% dos cursos.
De acordo com os dados apresentados, de todos os 12 cursos avaliados com conceito 5, três eram vinculados a Universidades.
Apenas quatro cursos em Universidades ficaram sem conceito ( SC ).
Esse tipo de Organização Acadêmica teve o conceito 3 como modal, com 18 cursos ( 32,1 \% ).
Os demais cursos de Universidades avaliados receberam os conceitos 1 ( três cursos ), 2 ( 16 cursos ), 4 ( 12 cursos ) e 5 ( três cursos, como já mencionado ).
Entre os cursos em Centros Universitários, o conceito modal também foi 3, com 15 cursos ( 57,7 \% ).
Neste tipo de organização acadêmica apenas dois cursos ficaram sem conceito ( SC ), enquanto nenhum recebeu os conceitos 1 ou 5.
Os outros cursos neste tipo de Organização Acadêmica receberam os conceitos 2 ( seis cursos ) e conceito 4 ( três cursos ).
Nas Faculdades, apenas um dos 96 cursos ficou sem conceito ( SC ).
O conceito modal ficou entre os conceitos 2 e 3, cada um com 33 cursos ( 34,4 \% ).
Dos demais cursos neste tipo de Organização Acadêmica, seis receberam conceito 1, 14 conceito 4 e nove cursos, o conceito 5.
Considerando -se separadamente as regiões brasileiras, verifica -se que, na região Norte, as Universidades não participaram com nenhum curso assim como os Centros Universitários.
As Faculdades concentraram todos os sete cursos participantes na região Norte.
Nenhum ficou sem conceito ( SC ) ou recebeu os conceitos 1 ou 5.
Os demais receberam conceito 2 ( cinco cursos, conceito modal ), 3 ( um curso ) e 4 ( também um curso ).
Na região Nordeste, as Universidades participaram com somente um curso na Área de Tecnologia em Processos Gerenciais da região, o qual foi avaliado com o conceito 1.
Os Centros Universitários também contaram com somente um curso, o qual foi avaliado com o conceito 2.
As Faculdades foram representadas por 18 cursos na região Nordeste.
O conceito modal foi 2, com seis cursos.
Dos cursos restantes, três receberam conceito 1, cinco conceito 3, dois conceito 4 e outros dois cursos, conceito 5.
Na região Sudeste, as Universidades concentraram 31 dos 89 cursos da região.
Entre os cursos em Universidades na região, o conceito modal ficou entre os conceitos 2 e 3, com nove cursos cada, e apenas dois cursos ficaram sem conceito ( SC ).
Os demais cursos receberam os conceitos 1 ( dois cursos ), 4 ( sete cursos ) e 5 ( outros dois cursos ).
Os Centros Universitários participaram com 17 cursos na região Sudeste, dos quais dez obtiveram conceito modal, 3, e apenas dois ficaram sem conceito ( SC ).
Os demais receberam o conceito 2 ( três cursos ) e conceito 4 ( dois cursos ).
Nenhum dos cursos nesta combinação de organização acadêmica e região recebeu conceitos 1 ou 5.
As Faculdades foram representadas por 41 cursos na região Sudeste, que se distribuíram nos conceitos 1 ( um curso ), 2 ( 13 cursos ), 3 ( 14 cursos, conceito modal ), 4 ( sete cursos ) e 5 ( cinco cursos ), enquanto um único curso ficou sem conceito ( SC ).
Dos 50 cursos da região Sul, 20 eram de Universidades, para os quais o conceito modal foi 3, com oito cursos.
Nesse tipo de organização, apenas um dos cursos ficou sem conceito ( SC ) e os demais receberam os conceitos 2 ( seis cursos ), 4 ( quatro cursos ) e 5 ( um curso ).
Os Centros Universitários da região Sul tiveram cinco dos seis cursos participantes no conceito modal, 3.
O outro curso recebeu conceito 2.
Foram 24 os cursos vinculados a Faculdades na região Sul, e nenhum destes ficou sem conceito ( SC ).
Sete desses cursos receberam conceito 2, 12 cursos receberam conceito 3, o conceito modal, três cursos receberam conceito 4 e os dois outros, conceito 5.
Na região Centro-Oeste, quatro dos 12 cursos eram de Universidades.
Nesse tipo de organização, apenas um curso ficou sem conceito ( SC ), enquanto os outros três cursos foram avaliados com os conceitos 2, 3 e 4 ( um curso cada ).
Os Centros Universitários da região Centro-Oeste contaram com dois cursos, dos quais um recebeu o conceito 2 e o outro, o conceito 4.
Nenhum curso ficou sem conceito ( SC ).
Dos seis cursos em Faculdades na região Centro-Oeste, nenhum ficou sem conceito ( SC ), dois receberam conceito 1, outros dois conceito 2, um conceito 3 e o outro conceito 4.

% matched lemmas: alocar
\end{textsample}
